\section{Protocol experimental}\label{sec:protocol}

Totes les execucions tenen el pressupost que fixa l'enunciat, $T = 300.000$ transicions, declarat abans d'executar res. DDPG i TD3 complet es comparen amb una llavor (la 0), com demana l'enunciat. A l'ablació, cada variant s'entrena amb les llavors 0--9 en les dues configuracions: l'enunciat en demana de 3 a 5, però la variància entre llavors és alta i vam preferir-ne deu. TD3 complet desa el \emph{checkpoint} de $T/3$ que demana l'enunciat, però totes les variants s'entrenen des de zero.

Cada 1.000 transicions s'avalua la política sense soroll en 5 episodis amb les mateixes llavors d'entorn (10000--10004) a totes les avaluacions. Abans de veure cap resultat vam fixar dues mètriques: la mitjana de les últimes 10 avaluacions (L10, 50 episodis) i l'àrea sota la corba d'avaluació dividida per $T$ (AUC/T). Les diferències respecte de TD3 complet de la mateixa configuració s'aparellen per llavor, amb un interval de confiança (IC) del 95~\% per \emph{bootstrap} (10.000 remostrejos). Com a anàlisi complementària, no fixada d'entrada, s'hi afegeix un test de permutació aparellat exacte amb correcció de Holm per a les tres comparacions de cada configuració. També es mesuren la fracció d'episodis amb retorn $\geq 200$, el llindar amb què es considera resolt l'entorn \parencite{gymnasium}, i la fracció d'avaluacions amb els cinc episodis truncats. A la configuració B, cada execució desa també les xarxes finals i les avalua en 100 episodis amb llavors noves (20000--20099).

El biaix del crític es mesura com a $Q_1(s_t,\mu(s_t)) - G_t$, on $G_t = \sum_k \gamma^k r_{t+k}$ és el retorn descomptat observat en els episodis d'avaluació. En els episodis truncats s'exclouen els estats amb menys de 300 passos per davant, perquè el seu $G_t$ queda incomplet. Les execucions es van fer en CPU, en dues màquines virtuals Azure Standard\_D4s\_v4 repartides per llavors i no per variants.
